% Propietario conservado: /Users/ruben/Documents/ChatGPT/jueces y controles/REVISION_ARGUMENTAL_APERTURAS_CIERRES_20260915/II/manuscript_es/sections/09c_entropia_sectorial.tex
\subsection{Ecuación de estado de la información sectorial}
\label{ii:sec:ii-ecuacion-informacion}

La incidencia ya construida proporciona dieciocho factores tríticos en
cada sector. Esta multiplicidad permite escribir explícitamente la
relación entre el área, la distribución de ocupaciones y la entropía de
la purificación anterior. Conservamos el estado reducido allí definido
y abreviamos $\Delta=\Delta_{\rm mod}$. Para $m\in\{90,120\}$ ponemos
\begin{equation}
 x_m=e^{-m\Delta},\qquad z_m=1+x_m+x_m^2,\qquad
 f_m=\frac{x_m+2x_m^2}{z_m},\qquad
 v_m=\frac{x_m+4x_m^2}{z_m}-f_m^2.
 \label{ii:eq:ii-sector-momentos}
\end{equation}
Aquí $f_m$ y $v_m$ son, respectivamente, la media y la varianza de
la ocupación de un factor del sector $m$. El parámetro $\Delta$ conserva
el estatuto y la normalización del estado reducido previamente
especificado.

\subsubsection{Purificación cuántica bipartita de los sectores de incidencia}
\label{ii:sec:ii-area-purificacion-producto}

La descomposición de Schmidt anterior se concreta en cada enlace.
Para un peso real finito \(w\), definimos
\[
 z(w)=1+e^{-w}+e^{-2w},\qquad
 \rho(w)=\frac1{z(w)}\sum_{n=0}^2e^{-wn}|n\rangle\langle n|.
\]
Los pesos sectoriales son \(w_m=m\Delta\), con \(m=90,120\);
por tanto, \(z(w_m)=z_m\).

\begin{proposition}[Purificación termocampo doble de la incidencia sectorial]
\label{ii:prop:ii-area-purificacion-producto}
El vector normalizado
\begin{equation}
 |\operatorname{TFD}(w)\rangle
 =\frac1{\sqrt{z(w)}}\sum_{n=0}^2
       e^{-wn/2}|n\rangle_A\otimes|n\rangle_{\bar A}
 \label{ii:eq:ii-area-tfd-enlace}
\end{equation}
purifica \(\rho(w)\). El producto
\begin{equation}
 |\Psi_\Delta\rangle
 =|\operatorname{TFD}(90\Delta)\rangle^{\otimes18}
  \otimes|\operatorname{TFD}(120\Delta)\rangle^{\otimes18}
 \label{ii:eq:ii-area-tfd-producto}
\end{equation}
tiene reducción \(\rho_A(\Delta)\) sobre los treinta y seis enlaces
originales. Su entropía de entrelazamiento para la bipartición
\(A\mid\bar A\) es \(s(\Delta)\).
\end{proposition}
\begin{proof}
La norma al cuadrado es \(z(w)^{-1}\sum_{n=0}^2e^{-wn}=1\).
En la traza parcial, la ortogonalidad del segundo factor elimina los
términos de índices distintos; los restantes forman \(\rho(w)\).
La traza parcial conmuta con el producto tensorial. Se obtiene así
\(\rho(90\Delta)^{\otimes18}\otimes\rho(120\Delta)^{\otimes18}\),
que coincide con el estado normalizado definido anteriormente.
La igualdad entre la entropía de una reducción y el entrelazamiento de
esta purificación se lee directamente de sus coeficientes de Schmidt.
\end{proof}

Los dos factores de la bipartición quedan determinados dentro de la
construcción. La realización física del segundo factor conserva esa
identificación de subsistema y su álgebra de observables.

\subsubsection{Momentos de ocupación y contenido informacional}

\begin{proposition}[Área media y entropía de la purificación sectorial]
En esta familia finita de estados,
\begin{align}
 \langle N\rangle_\Delta&=18(f_{90}+f_{120}),\\
 \langle D\rangle_\Delta&=18(90f_{90}+120f_{120}),\\
 \frac{\langle\widehat{\mathcal A}\rangle_\Delta}{\ell_P^2}
 &=18\gamma_{\HMT}a_0(f_{90}+f_{120}),
 \label{ii:eq:ii-area-media}\\
 s(\Delta)&=18\log z_{90}+18\log z_{120}
                  +\Delta\langle D\rangle_\Delta.
 \label{ii:eq:ii-entropia-sectorial}
\end{align}
La entropía dimensional es $S(\Delta)=k_Bs(\Delta)$, con el
transductor térmico de la sección~\ref{ii:sec:ii-boltzmann}.
\end{proposition}
\begin{proof}
En cada factor los pesos normalizados son
$(1,x_m,x_m^2)/z_m$. Su primer momento da $f_m$ y su segundo
momento da la primera fracción de $v_m$. La factorización tensorial
y la aditividad de los operadores número producen las tres primeras
identidades. Finalmente,
$\log\rho_A=-(\log Z)I-\Delta D$, con
$Z=z_{90}^{18}z_{120}^{18}$. Tomar la traza contra $-\rho_A$
da~\eqref{ii:eq:ii-entropia-sectorial}. La purificación especificada
anteriormente tiene esta misma entropía de entrelazamiento.
\end{proof}

La pérdida de uniformidad tiene una lectura cuantitativa propia.
En dimensión \(d_\partial=3^{36}\), sea
\(\tau_\partial=I/d_\partial\). Definimos
\begin{equation}
 I_{\rm or}(\Delta)=36\log3-s(\Delta)
 =D\bigl(\rho_A(\Delta)\Vert\tau_\partial\bigr)\ge0.
 \label{ii:eq:ii-area-deficit-informacional}
\end{equation}
La igualdad se obtiene sustituyendo
\(\log\tau_\partial=-36\log3\,I\). Si \(p_\nu\) son los
autovalores de \(\rho_A\), la convexidad de \(t\log t\) da
\(\sum_\nu p_\nu\log(d_\partial p_\nu)\ge0\), con igualdad
exactamente para el estado uniforme. Para esta familia, ello ocurre en
\(\Delta=0\). El déficit es una información adimensional; el cuanto
areal \(\gamma_{\HMT}a_0\) es otra lectura. Cada expresión conserva
su definición aunque sus evaluaciones puedan estar próximas.

\begin{proposition}[Respuesta informacional y fluctuaciones de borde]
Para todo $\Delta$ real finito,
\begin{align}
 \frac{\dd\langle D\rangle_\Delta}{\dd\Delta}
 &=-\operatorname{Var}_\Delta(D)
   =-18(90^2v_{90}+120^2v_{120}),\\
 \frac{\dd s}{\dd\Delta}
 &=-\Delta\operatorname{Var}_\Delta(D),
 \label{ii:eq:ii-respuesta-entropica}\\
 \frac{\dd\langle\widehat{\mathcal A}\rangle_\Delta}{\dd\Delta}
 &=-18\gamma_{\HMT}a_0\ell_P^2(90v_{90}+120v_{120}).
 \label{ii:eq:ii-respuesta-areal}
\end{align}
\end{proposition}
\begin{proof}
La suma que define $Z$ es finita. Su derivación da
$\partial_\Delta\log Z=-\langle D\rangle_\Delta$ y
$\partial_\Delta^2\log Z=\operatorname{Var}_\Delta(D)$.
Los sectores y sus factores son independientes en el estado producto;
por ello se suman sus varianzas. Al derivar
\eqref{ii:eq:ii-entropia-sectorial}, los términos
$-\langle D\rangle_\Delta$ y $\langle D\rangle_\Delta$
se cancelan. Para el área se utiliza
$\partial_\Delta f_m=-mv_m$ en~\eqref{ii:eq:ii-area-media}.
\end{proof}

Estas ecuaciones relacionan tres niveles de una misma construcción.
La incidencia fija las multiplicidades y los pesos sectoriales; el
estado reducido fija sus probabilidades; las fluctuaciones de esas
ocupaciones determinan la respuesta del área y de la entropía.
La escala $\gamma_{\HMT}a_0\ell_P^2$ convierte la ocupación en área,
mientras $k_B$ convierte la información en entropía dimensional.
En la sección gravitatoria, $\ell_P^2=G\hbar/c^3$ explicita
la dependencia de esa unidad areal respecto de acción y gravitación.
La comparación con un horizonte conserva adicionalmente el mapa
geométrico y el estado que se especifican allí.

\subsection{Inversión de ocupaciones e invariantes entrópicos}
\label{ii:sec:ii-inversion-sectorial}

Sea $\mathcal R$ la involución unitaria que intercambia $|0\rangle$
y $|2\rangle$ y fija $|1\rangle$ en cada uno de los treinta y seis
factores tríticos. Su acción sobre las ocupaciones conserva toda la
configuración microscópica mediante un cambio biyectivo de sus etiquetas.

\begin{theorem}[Entropía conservada y área complementaria]
La involución anterior satisface
\begin{align}
 \mathcal R N\mathcal R^*&=72I-N,&
 \mathcal R D\mathcal R^*&=7560I-D,\\
 \rho_A(-\Delta)&=\mathcal R\rho_A(\Delta)\mathcal R^*,&
 s(-\Delta)&=s(\Delta),
 \label{ii:eq:ii-inversion-entropia}\\
 \langle\widehat{\mathcal A}\rangle_{-\Delta}
 +\langle\widehat{\mathcal A}\rangle_\Delta
 &=72\gamma_{\HMT}a_0\ell_P^2.&&
 \label{ii:eq:ii-area-complementaria}
\end{align}
La entropía relativa entre ambas orientaciones es
\begin{equation}
 D\bigl(\rho_A(\Delta)\Vert\rho_A(-\Delta)\bigr)
 =\Delta\bigl(7560-2\langle D\rangle_\Delta\bigr).
 \label{ii:eq:ii-distancia-orientaciones}
\end{equation}
\end{theorem}
\begin{proof}
Cada ocupación local $n$ se transforma en $2-n$. Por tanto,
$N_m\mapsto36I-N_m$ y
$D\mapsto36(90+120)I-D=7560I-D$.
Además, $z_m(-\Delta)=e^{2m\Delta}z_m(\Delta)$,
de donde $Z(-\Delta)=e^{7560\Delta}Z(\Delta)$.
Sustituir ambas identidades en el estado normalizado prueba la
conjugación unitaria. La invariancia espectral prueba la igualdad
de entropías; la transformación de $N$ prueba la suma de áreas.
Finalmente, al restar los dos logaritmos de los estados y tomar
su esperanza en $\rho_A(\Delta)$ se obtiene
\eqref{ii:eq:ii-distancia-orientaciones}.
\end{proof}

En $\Delta=0$ el estado reducido es uniforme y
$s(0)=36\log3$. En el límite $\Delta\to+\infty$ sobrevive
la configuración de ocupación nula; la entropía tiende a cero.
La inversión lleva ese límite a la configuración de ocupación máxima.
En ambos extremos la entropía es la misma y el área distingue las
configuraciones. Para $\Delta>0$, la respuesta
\eqref{ii:eq:ii-respuesta-entropica} es estrictamente negativa porque
el estado tiene rango completo y $D$ tiene más de un autovalor.
La ley conserva así la diferencia entre reversibilidad de la
configuración completa, igualdad de entropías y cambio de lectura areal.
Estas propiedades pertenecen a la familia sectorial especificada;
el transporte físico de la bipartición se trata en
la sección~\ref{ii:sec:confluencia}.

\subsection{Primera ley modular y variaciones finitas de área}
\label{ii:sec:ii-area-ley-modular}

Las derivadas de \eqref{ii:eq:ii-respuesta-entropica} varían el parámetro
\(\Delta\) de una familia de estados. La primera ley modular fija,
en cambio, un estado de referencia y considera variaciones del estado
observado alrededor de él. Tomemos
\(\rho_0=\rho_A(\Delta_0)\), con \(\Delta_0\) real y finito,
y mantengamos fijos \(\Delta_0,\gamma_{\HMT},a_0,\ell_P\).
Se conserva la rama \(\gamma_{\HMT}>0\), \(a_0>0\).
Las áreas sectoriales normalizadas son
\begin{equation}
 \mathcal A_m^{\rm n}=\gamma_{\HMT}a_0N_m,
 \qquad
 \widehat{\mathcal A}_m=\ell_P^2\mathcal A_m^{\rm n},
 \qquad m\in\{90,120\}.
 \label{ii:eq:ii-area-sector-normalizado}
\end{equation}
El logaritmo del estado de referencia adopta la forma
\begin{equation}
 K_0=-\log\rho_0
 =c_0I+\beta_{90}\mathcal A_{90}^{\rm n}
         +\beta_{120}\mathcal A_{120}^{\rm n},
 \qquad
 \beta_m=\frac{m\Delta_0}{\gamma_{\HMT}a_0},
 \quad c_0=\log Z(\Delta_0).
 \label{ii:eq:ii-area-beta-sectorial}
\end{equation}
En particular, \(\beta_{120}=\tfrac43\beta_{90}\); para
\(\Delta_0\ne0\), su cociente es \(4/3\). Los coeficientes
surgen de expresar el mismo operador modular en unidades areales.

\begin{theorem}[Primera ley modular en coordenadas areales]
\label{ii:thm:ii-area-primera-ley}
Sea \(\sigma(\varepsilon)\) una familia diferenciable de estados
normalizados en un intervalo que contiene cero, con
\(\sigma(0)=\rho_0\). Entonces
\begin{equation}
 \left.\frac{\dd s(\sigma(\varepsilon))}{\dd\varepsilon}\right|_0
 =\sum_{m=90,120}\beta_m
   \left.\frac{\dd\langle\mathcal A_m^{\rm n}\rangle_
                       {\sigma(\varepsilon)}}{\dd\varepsilon}\right|_0.
 \label{ii:eq:ii-area-primera-ley}
\end{equation}
\end{theorem}
\begin{proof}
El estado \(\rho_0\) tiene rango completo. En su entorno, la derivada
de la traza de \(-\sigma\log\sigma\) es
\(-\Tr[\dot\sigma(\log\rho_0+I)]\); se obtiene diagonalizando
\(\rho_0\) en la fórmula de la derivada espectral de la traza.
La normalización implica \(\Tr\dot\sigma=0\), de modo que la
derivada es \(\Tr(\dot\sigma K_0)\). Sustituir
\eqref{ii:eq:ii-area-beta-sectorial} cancela también \(c_0I\) y
produce la identidad.
\end{proof}

\begin{theorem}[Identidad finita de área y déficit de entropía relativa]
\label{ii:thm:ii-area-ley-finita}
Para cualquier estado normalizado \(\sigma\) del mismo espacio finito,
definamos \(\Delta_\sigma\langle B\rangle
=\Tr[(\sigma-\rho_0)B]\). Se cumple
\begin{equation}
 \boxed{
 s(\sigma)-s(\rho_0)
 =\sum_{m=90,120}\beta_m
      \Delta_\sigma\langle\mathcal A_m^{\rm n}\rangle
   -D(\sigma\Vert\rho_0).}
 \label{ii:eq:ii-area-ley-finita}
\end{equation}
Aquí \(D(\sigma\Vert\rho_0)=\Tr[\sigma(\log\sigma-\log\rho_0)]\)
es finita, no negativa, y se interpreta con \(0\log0=0\).
\end{theorem}
\begin{proof}
Por definición,
\(D(\sigma\Vert\rho_0)=-s(\sigma)+\Tr(\sigma K_0)\), mientras
\(s(\rho_0)=\Tr(\rho_0K_0)\). Restar da
\(s(\sigma)-s(\rho_0)=\Delta_\sigma\langle K_0\rangle
-D(\sigma\Vert\rho_0)\). Sustituir
\eqref{ii:eq:ii-area-beta-sectorial} demuestra la fórmula porque las
trazas de ambos estados son uno.

Para ver la positividad, diagonalicemos \(\sigma\) con autovalores
\(p_i\) y vectores \(u_i\), y \(\rho_0\) con autovalores
\(q_j>0\) y vectores \(v_j\). Pongamos
\(r_i=\langle u_i,\rho_0u_i\rangle\). La concavidad del logaritmo
proporciona
\(\sum_j|\langle u_i,v_j\rangle|^2\log q_j\le\log r_i\).
De aquí,
\(D(\sigma\Vert\rho_0)\ge\sum_i p_i\log(p_i/r_i)\ge0\);
la última desigualdad es la convexidad de \(t\log t\) aplicada
con pesos \(r_i\), que suman uno. La positividad estricta de los
\(q_j\) garantiza que todos los logaritmos de referencia son finitos.
\end{proof}

La identidad finita conserva la distancia informacional respecto del
estado de referencia; su término tangente es
\eqref{ii:eq:ii-area-primera-ley}. Para las áreas físicas, los
coeficientes correspondientes son \(\beta_m/\ell_P^2\).
Multiplicar la ecuación por el transductor común \(k_B\) publica
la entropía dimensional. Así, la relación entre geometría e información
permanece exacta también para variaciones finitas, con su corrección
explícita y sin modificar los parámetros del estado de referencia.
